\section{Completion in the canonical coordinate}
\label{comp:section}

The local passage comparison allows endpoints to be chosen within
small ports.  We now identify distances between arbitrary vertices,
including vertices approaching the exploration root.  The intrinsic
GH input supplies a compact metric sphere; the task is to identify
its points with their limiting conformal locations.

Fix $q\in(0,4)$.  We use the deterministic normalization $a_n$ in
Theorem~\ref{intro:intrinsic-input}, with the continuum metric
normalization specified there.  Put
\[
 V_n=V(M_n),\qquad \delta_n=a_n d_n,
 \qquad z_n(v)=\phi_n(v),\qquad
 \mu_n=(\phi_n)_*m_n.
\]
The metric $d_n$ here uses every original primal edge.  In particular,
shortest paths used below are paths in the original primal graph.

We keep two coordinates distinct.  The original oriented peanosphere
coordinate has field $h^{\mathrm o}$, metric $D^{\mathrm o}$, area
$\mu^{\mathrm o}$ and exploration root $p_*^{\mathrm o}=\infty$.
It has the quantum-sphere and independent imaginary-field laws used
in Sections~\ref{enc:section} and~\ref{ker:section}.  Write
\[
 x_n(v)=H_n(v),\qquad
 f_n=\phi_n\circ H_n^{-1}\longrightarrow f
\]
along a joint extraction of Theorem~\ref{conf:canonical}.  The limit
$f$ may be M\"obius or anti-M\"obius, and its orientation may depend
on all the retained data.  Define the canonical geometric objects by
\begin{equation}
 \label{comp:two-coordinates}
 \mu=f_*\mu^{\mathrm o},\qquad p_*=f(p_*^{\mathrm o}),\qquad
 D(z,w)=D^{\mathrm o}(f^{-1}(z),f^{-1}(w)).
\end{equation}
These definitions transport geometry sample by sample.  Field-density
and local-kernel arguments use the oriented peanosphere field and,
when constructed below, its orientation-preserving reembedding by
fresh area marks.  We verify the field law in that reembedding
explicitly; no field-law assertion is inferred from the possibly
adaptive orientation of $f$.  Port and circuit conclusions in the
canonical coordinate are the
pathwise transports of the already proved conclusions for $x_n$;
uniform convergence of $f_n$ and $f_n^{-1}$ preserves their strict
margins.

Let $T$ be the orientation-preserving M\"obius map that takes the
three actual limiting normalization samples in the original coordinate
to $0,1,\infty$, and set
\[
 \mu^+=T_*\mu^{\mathrm o},\qquad
 D^+(z,w)=D^{\mathrm o}(T^{-1}(z),T^{-1}(w)).
\]
This pair has the prescribed canonically marked quantum-sphere law.
Write $\mathfrak r(z)=\overline z$ on $\Chat$, with
$\mathfrak r(\infty)=\infty$.  It acts on a measure by pushforward,
on a distance graph by
$(z,w,r)\mapsto(\mathfrak r z,\mathfrak r w,r)$, and on a loop
collection by composing every loop with $\mathfrak r$.  These actions
are continuous in the topologies of the main theorem; the graph and
loop distances are unchanged.  The map $f$ is either $T$ or
$\mathfrak r\circ T$.  We will remove this common reflection at
the level of the joint law after identifying the metric graph.

\begin{theorem}[Completion of the metric graph]
\label{comp:metric-graph}
In every joint subsequential extraction supplied by
Theorems~\ref{conf:canonical} and~\ref{met:normalizers}, one may further
extract the intrinsic spaces $(V_n,\delta_n)$ and their spatial distance
graphs.  Almost surely, every resulting limit of the latter is
\begin{equation}
 \label{comp:target-graph}
 \mathcal G_D
   =\{(z,w,D(z,w)):z,w\in\Chat\}.
\end{equation}
Consequently
\[
 (\mu_n,K_n(a_n))\ \Longrightarrow\ (\mu^+,\mathcal G_{D^+})
\]
through all positive integers.  In any convergent joint extraction
with additional reference or local observations, its metric
component is the same graph \eqref{comp:target-graph}.  The
Hausdorff metric for the second component is that induced by
$\rho(z,z')+\rho(w,w')+|r-r'|$.
\end{theorem}

The proof compares the passage scale $B(n)$ of
Theorem~\ref{met:normalizers} with the intrinsic scale $a_n$.
We first show that $a_nB(n)$ is bounded and retain a deterministic
subsequential limit $\alpha\geq0$; let $\lambda$ be the passage
comparison constant on the same extraction.  There are two possibilities.
If $\alpha>0$, the lower passage comparison and small primal
separators identify the location relation as a homeomorphism, with
metric factor $\alpha\lambda$.  If $\alpha=0$, those separators
force every intrinsic point except one to have the same spatial
location.  Two separated local escape tests will rule out that
concentration.  Once $\alpha>0$, the intrinsic diameter law fixes
$\alpha\lambda=1$.

We begin by retaining the full location relation.  The intrinsic
input identifies each marginal limit of $(V_n,\delta_n)$ as a
nontrivial compact metric sphere, but gives no identification with
the coordinate $z_n$.

\subsection{Location relations and separators}

\begin{lemma}[Retaining every vertex]
\label{comp:relation}
Suppose that, on a coupled subsequence,
$(V_n,\delta_n)$ converges in Gromov--Hausdorff distance to a compact
metric space $(Y,d_Y)$, and that $z_n(V_n)$ becomes dense in $\Chat$.
There are maps $j_n:V_n\to Y$ and numbers $e_n\downarrow0$ such that
\begin{equation}
 \label{comp:approximation}
 \left|d_Y(j_n(u),j_n(v))-\delta_n(u,v)\right|\le e_n,
 \qquad
 \sup_{y\in Y}\dist_{d_Y}(y,j_n(V_n))\le e_n.
\end{equation}
After a further subsequence, the finite sets
\[
 R_n=\{(j_n(v),z_n(v)):v\in V_n\}
\]
converge in Hausdorff distance in $Y\times\Chat$ to a compact relation
$R$ that projects onto both factors.  On the same subsequence,
\begin{equation}
 \label{comp:relation-graph}
 K_n(a_n)\longrightarrow
 \mathcal K(R):=
 \{(z,w,d_Y(y,y')):(y,z),(y',w)\in R\}.
\end{equation}
\end{lemma}

\begin{proof}
A correspondence of distortion tending to zero in a
Gromov--Hausdorff realization gives \eqref{comp:approximation}: choose
one related point of $Y$ for each vertex.  The image is a net with
error tending to zero.  The space of nonempty compact subsets of the
compact space $Y\times\Chat$ is compact, so a subsequence of $R_n$
converges.  Net density in the two factors proves the two surjectivity
assertions.

For the upper inclusion in \eqref{comp:relation-graph}, take any
converging sequence of triples from $K_n(a_n)$ and extract the two
$Y$-coordinates of its vertices.  They give two pairs in $R$, and
\eqref{comp:approximation} identifies the third coordinate.  For the
reverse inclusion, approximate any two pairs of $R$ by pairs of $R_n$
at the same indices.  Their distances again converge by
\eqref{comp:approximation}.  This proves both Hausdorff inclusions.
\end{proof}

The mesh assertion of Corollary~\ref{conf:mesh-area} supplies the
spatial density needed in this lemma: every flag contains an original
vertex, and the maximum spherical flag diameter tends to zero.  The
maps $j_n$ need not be continuous or agree with the spatial coordinate.

Here is the elementary use of the topology of the intrinsic limit.
For a set of vertices $S_n$, write $G_n\setminus S_n$ for the graph
obtained by deleting those vertices and their incident edges.

\begin{lemma}[A collapsing separator]
\label{comp:avoid-separator}
Assume \eqref{comp:approximation}.  Let $y_0\in Y$, and suppose that
$Y\setminus\{y_0\}$ is path connected.  If
\[
 \sup_{v\in S_n}d_Y(j_n(v),y_0)\longrightarrow0,
\]
then, for any vertices $u_n,v_n$ with
$j_n(u_n)\to y$ and $j_n(v_n)\to y'$ in
$Y\setminus\{y_0\}$, the two vertices are in the same component of
$G_n\setminus S_n$ for all sufficiently large $n$.
The conclusion also holds when $y=y'$.
\end{lemma}

\begin{proof}
Choose a continuous path $P$ from $y$ to $y'$ in
$Y\setminus\{y_0\}$.  If the endpoints coincide, take the constant
path.  Its compact image has distance $b>0$ from $y_0$.  By uniform
continuity, choose a finite partition of $P$ into pieces whose
endpoints are at $d_Y$-distance less than $b/10$.  Approximate the
partition points by $j_n(V_n)$, using $u_n,v_n$ at the ends.  For large
$n$, consecutive chosen vertices have $\delta_n$-distance less than
$b/5$, and each chosen vertex has $Y$-distance at least $4b/5$ from
$y_0$.

Join consecutive vertices by graph shortest paths.  Every vertex of
such a path is at $\delta_n$-distance less than $b/5$ from its initial
vertex, so \eqref{comp:approximation} keeps its $Y$-image at distance
at least $b/2$ from $y_0$.  It therefore avoids $S_n$.  Concatenation
proves the assertion.  The same proof shows the following useful
version: for this fixed path $P$, it is enough that all of $j_n(S_n)$
eventually lies within distance $b/4$ of $y_0$.
\end{proof}

Every point of a topological sphere has a path-connected complement.
Thus this lemma applies to the intrinsic limit furnished by
Theorem~\ref{intro:intrinsic-input}.  When $S_n$ is a primal circuit
or a connected primal skeleton, a path in $G_n\setminus S_n$ cannot
cross it on the embedded surface: primal edges meet only at their
incidence vertices.  The strict separation assertions in
Corollary~\ref{top:circuits} therefore apply to these graph paths,
including in maps with loops or multiple edges.

\subsection{The comparison scale and the exploration root}

\begin{lemma}[Small bypasses at an area mark]
\label{comp:root-bypass}
Almost surely there are rectifiable closed paths $L_k$ in
$\Chat\setminus\{p_*\}$, each separating $p_*$ from the complement of
a neighborhood of $p_*$, whose spatial diameters and $D$-lengths tend
to zero.  Consequently, for $z,w\ne p_*$, the internal distance in
$\Chat\setminus\{p_*\}$ equals $D(z,w)$.  The metric completion of
that punctured space is $(\Chat,D)$.
\end{lemma}

\begin{proof}
Work first with $D^{\mathrm o}$ in the original oriented sphere.
A fixed orientation-preserving inversion puts
$p_*^{\mathrm o}=\infty$ at zero.  By
Proposition~\ref{enc:sphere-representation}, near this mark the field
is represented by a GFF with a logarithmic insertion of strength
$\gamma$, a locally continuous term, and the finite random additive
constant that normalizes its area; the change to the unit-area law
has a positive integrable density.  We may first work with
\[
 h^\gamma=h^0-\gamma\log|\cdot|,
\]
where $h^0$ is a whole-plane GFF normalized on the unit circle.
The local continuous term changes lengths by factors bounded above
and below on a fixed smaller disk, by Weyl scaling.

For $0<r<1/4$, consider the annulus $r<|z|<2r$.  The
around-annulus estimate for $h^\gamma$ gives a disconnecting closed
path of length at most
\begin{equation}
 \label{comp:root-annulus-bound}
 r^{\xi(Q-\gamma)}e^{\xi h^0_r(0)}T_r,
\end{equation}
where the nonnegative variables $T_r$ form a tight family.  This is
the around-distance estimate, not just a point-to-circle estimate;
see Definition~3.7 and the proof of Proposition~3.14 in
\cite{DubedatEtAlWeakMetric}.  It also follows by scaling a unit annulus, using
the spatial coordinate rule and the tight unit-annulus internal
distances.  The normalized field
$h^0(r\cdot)-h^0_r(0)$ has the same annular law at every $r$.

Now
\[
 Q-\gamma=\frac{4-\gamma^2}{2\gamma}>0,
 \qquad
 \frac{h^0_r(0)}{\log(1/r)}\longrightarrow0
 \quad\hbox{in probability}.
\]
Indeed the numerator is centered Gaussian of variance
$\log(1/r)$.  Thus \eqref{comp:root-annulus-bound} tends to zero in
probability; no independence between its factors is needed.  Bounded
local continuous perturbations, the finite area-normalizing factor,
and an absolutely continuous probability weighting preserve this
convergence.  Choose deterministic radii $r_k\downarrow0$ so that
the probabilities of a length exceeding $2^{-k}$ are summable.
Borel--Cantelli gives such paths almost surely in the original
oriented sphere.  Their diameters tend to zero because they stay
in the corresponding shrinking annuli.  Transport them by $f$.
Equation~\eqref{comp:two-coordinates} preserves their metric lengths,
and continuity of $f$ gives vanishing canonical spatial diameter.
This obtains the asserted $L_k$ without changing any field law.

Let $z,w\ne p_*$.  Take a $D$-geodesic between them.  If it passes
through $p_*$, replace the portion between its first and last
intersections with $L_k$ by a portion of $L_k$; both endpoints lie
outside the annulus for large $k$.  The replacement avoids $p_*$ and
adds at most the length of $L_k$.  A closed path with nonzero winding
can, for this purpose, be replaced by a separating simple subcircuit
of its trace; equivalently one can use the first and last crossing
of its separating trace and either connecting portion of the closed
path.  The internal distance is therefore at most $D(z,w)$ after
letting $k\to\infty$.  The reverse inequality is immediate.
Finally the punctured sphere is $D$-dense by the topology of the LQG
metric, and $(\Chat,D)$ is complete.  This identifies its completion.
\end{proof}

\begin{lemma}[Bounded deterministic comparison scales]
\label{comp:scales}
The sequence $a_nB(n)$ is bounded.  On every deterministic subsequence
one may therefore pass to a further deterministic subsequence on
which
\begin{equation}
 \label{comp:alpha}
 a_nB(n)\longrightarrow\alpha\in[0,\infty).
\end{equation}
In a joint extraction on this subsequence, let $\lambda$ be the
deterministic positive multiplier of
Theorem~\ref{met:normalizers}.  In $\delta_n$-units, local port
following has upper factor $\alpha\lambda$.  If $\alpha>0$, the
lower certificate comparison has that same factor.

Almost surely, about every spatial point there are shrinking strict
primal separating circuits whose limiting $\delta_n$-diameters tend
to zero.  Moreover, if $\alpha=0$, then for every $\varepsilon>0$
there are connected primal skeletons of $\delta_n$-diameter tending
to zero whose complementary spatial components have diameter at
most $\varepsilon$ for all sufficiently large $n$.
\end{lemma}

\begin{proof}
Choose a fixed strictly buffered annular test in an ordinary patch
of the original oriented coordinate $x_n$, with vertices on both
sides.  The positive lower port comparison in
$B(n)$-units implies that, for some $b,\theta>0$ and all sufficiently
large $n$, every traversal of that annulus has cost at least $bB(n)$
on an event of probability at least $\theta$.  To obtain fixed
$b,\theta$, use the positive reference distance between its compact
boundary sets and the fixed positive lower bound for $\lambda$;
then truncate to a positive-probability event.  If the initial
choice of patch depends on the location of the exploration mark,
choose one member of the countable ordinary atlas with positive
probability and restrict to its strict containment event.

Every path between the two sides contains an annular traversal.
Hence the ambient graph diameter is at least $bB(n)$ on this event.
If $a_nB(n)$ were unbounded, a deterministic subsequence tending to
infinity would contradict the tightness of
$\diam(V_n,\delta_n)$ in Theorem~\ref{intro:intrinsic-input}.
This proves boundedness and \eqref{comp:alpha}.  Multiplying the
upper and lower comparisons by $a_nB(n)$ gives the stated factors.

Away from the original exploration root, the simultaneous circuit
conclusion is Lemma~\ref{met:vanishing-joins}, multiplied by the
bounded factors $a_nB(n)$.  Transporting the circuits by $f_n\to f$
gives the asserted conclusion away from $p_*$.  Its simultaneous
quantifier is relevant here: a point
chosen from a later metric relation is allowed.  The countable
strict port library and local internal H\"older bounds in that
lemma provide the assertion on every compact set off $p_*$.  At
$p_*$ use Lemma~\ref{comp:root-bypass}, cover each of its finitely
long reference paths by finitely many ordinary patches, and apply
port following and Corollaries~\ref{top:transverse} and~\ref{top:circuits}.
The resulting primal circuit has a strict central hole and lies in
a shrinking neighborhood of $p_*$.  Its limiting cost is at most
$\alpha\lambda$ times the reference length, with arbitrarily small
open losses.  These lengths tend to zero.

For the last assertion, choose finitely many strict central holes
covering the sphere, each bounded by one of the preceding reference
circuits and with filled side of spatial diameter less than
$\varepsilon$.  Join the circuits by finitely many reference paths
off $p_*$.  Such paths have finite $D$-length by
Lemma~\ref{comp:root-bypass}.  Use transverse ports at all required
intersections.  Lemma~\ref{top:separator-mesh} realizes this finite
construction as a connected primal skeleton.  Each complementary
component lies on the filled side of one of the covering circuits,
so has the stated diameter, allowing the fixed strict margins to
be chosen in advance.  The total cost in $B(n)$-units is bounded
along the extraction for this finite construction.  Multiplication
by $a_nB(n)\to0$ makes its $\delta_n$-diameter tend to zero.
\end{proof}

\subsection{Positive scale and continuity at all endpoints}

Fix a joint coupling on which the preceding convergences hold, and
retain the relation $R$ of Lemma~\ref{comp:relation}.  The intrinsic
limit $Y$ is almost surely homeomorphic to a sphere.  All remaining
topological arguments are deterministic on this event.

\begin{proposition}[Eliminating nontrivial fibers]
\label{comp:positive-fibers}
If $\alpha>0$, then $R$ is the graph of a homeomorphism
$\pi:Y\to\Chat$, and
\begin{equation}
 \label{comp:proportionality}
 d_Y(y,y')=\alpha\lambda D(\pi(y),\pi(y'))
 \qquad(y,y'\in Y).
\end{equation}
In particular,
\begin{equation}
 \label{comp:endpoint-modulus}
 \lim_{r\downarrow0}\limsup_{n\to\infty}
 \sup_{\substack{u,v\in V_n\\\rho(z_n(u),z_n(v))\le r}}
 \delta_n(u,v)=0.
\end{equation}
\end{proposition}

\begin{proof}
Write $c=\alpha\lambda>0$.  We first prove that every two related
pairs satisfy
\begin{equation}
 \label{comp:relation-lower}
 d_Y(y,y')\ge cD(z,z')
 \qquad\bigl((y,z),(y',z')\in R\bigr).
\end{equation}
Choose representing vertices and graph shortest paths between them.
Their $\delta_n$-lengths are bounded.  Parameterize them by this
length and keep them constant after their endpoints, on a common
bounded time interval.  Project the edges by the surface
homeomorphisms and the canonical coordinate.

For each $r>0$, any projected path of spherical diameter at least
$r$ makes a fixed positive advance in one member of a finite
ordinary annular cover.  Such a cover can avoid a ball of radius
much smaller than $r$ about $p_*$: a path with that diameter has a
subpath making a comparable advance outside the smaller ball.
The compact boundary sets in these finitely many tests have
positive $D$-separation.  The lower certificate comparison therefore
gives, for all sufficiently large $n$, a positive lower
$\delta_n$-cost for every such advance.  If this uniform lower
bound failed, a sequence of failing traversals would yield a
lower-certificate violation in one of the finitely many tests.
The maximum single-edge spatial diameter tends to zero and
$a_n\to0$, so the length-parameterized projections are
equicontinuous.  A subsequence converges uniformly to a continuous
path $P$ from $z$ to $z'$.

Every subarc of $P$ compactly off $p_*$ pays at least $c$ times the
$D$-distance of its ends.  Indeed a finite ordinary cover and
subdivision reduce it to the local lower comparisons, with costs
added over disjoint subpaths.  For $t>0$, remove the portion of $P$
between its first entry into and last exit from the $D$-ball of
radius $t$ about $p_*$.  The retained pieces are compactly off the
root and their endpoint distances sum to at least
$D(z,z')-2t$.  The same statement, with one retained piece, applies
if an endpoint is $p_*$.  If the ball is not visited there is no
removal.  Letting $t\downarrow0$ proves
\eqref{comp:relation-lower}.

Taking $y=y'$ in \eqref{comp:relation-lower} shows that $R$ assigns
each $y$ a unique spatial point.  Its graph is compact, so the
resulting surjection $\pi:Y\to\Chat$ is continuous.  In fact
\eqref{comp:relation-lower} gives its continuity directly in the
$D$-topology.

Suppose that a fiber over $z$ contains more than one point.  Take
the shrinking primal circuits about $z$ from
Lemma~\ref{comp:scales}.  For each fixed circuit test, extract its
$j_n$-image as a compact subset of $Y$.  Choose the tests so that
the diameters of these limiting subsets tend to zero.  After a
further subsequence of tests they converge to one point $y_0$.
Their spatial diameters tend to zero at $z$, so $\pi(y_0)=z$.
Choose $y\in\pi^{-1}(z)\setminus\{y_0\}$ and choose $x\in Y$
with $\pi(x)\ne z$.  Inequality~\eqref{comp:relation-lower} keeps
$x$ away from $y_0$.

There is a path in $Y\setminus\{y_0\}$ from $y$ to $x$ with
positive metric clearance from $y_0$.  Fix a sufficiently fine
circuit test that its limiting $Y$-image lies within one fourth of
this clearance from $y_0$.  Then take $n$ large.  Representatives
of $y$ lie in its strict central hole, and representatives of $x$
lie on its exterior side.  The last version of
Lemma~\ref{comp:avoid-separator} joins those representatives while
avoiding the entire circuit, contradicting its separation.
Thus the fibers are singletons.  A continuous bijection between
compact metric spaces is a homeomorphism.

To obtain \eqref{comp:endpoint-modulus}, a contrary sequence would
have spatially approaching pairs at distances bounded away from
zero.  Extract both $Y$-coordinates.  The relation then assigns
the same spatial point to two distinct $Y$-points, contradicting
the singleton-fiber conclusion.  This argument is uniform over
all vertices.

It remains to prove the upper bound in
\eqref{comp:proportionality}.  For $z,z'\ne p_*$ and any $e>0$,
Lemma~\ref{comp:root-bypass} supplies a rectifiable path compactly
in the punctured sphere, from $z$ to $z'$, of length at most
$D(z,z')+e$.  A finite cover by ordinary patches and the port upper
comparison give discrete following paths with arbitrarily small
spatial errors at their ends and limiting cost at most
$c(D(z,z')+e)$.  Equation~\eqref{comp:endpoint-modulus} now connects
their ends to any prescribed representing vertices at vanishing
additional cost: first make the fixed port errors small, then
take $n$ large, and finally let those errors tend to zero.
Let $e\downarrow0$.  Endpoints equal to $p_*$ follow by the
continuity of $D$ and of $\pi^{-1}$.  This proves the reverse of
\eqref{comp:relation-lower}, and hence
\eqref{comp:proportionality}.
\end{proof}

\subsection{The zero-scale alternative}

The topology of a sphere also describes what would happen if the
port scale vanished in GH units.  This description is stronger than
a statement about typical vertices.

\begin{proposition}[Concentration at one spatial point]
\label{comp:concentration}
If $\alpha=0$, there are unique points $y_0\in Y$ and $z_0\in\Chat$
such that
\begin{equation}
 \label{comp:cross-relation}
 R=(\{y_0\}\times\Chat)\ \cup\ (Y\times\{z_0\}).
\end{equation}
Moreover $z_0$ is a measurable function of the limiting distance
graph $\mathcal K(R)$ alone.  It is the unique point $z$ for which
\begin{equation}
 \label{comp:positive-diagonal}
 (z,z,r)\in\mathcal K(R)\quad\hbox{for some }r>0.
\end{equation}
\end{proposition}

\begin{proof}
For each $k\ge1$, take a connected skeleton from
Lemma~\ref{comp:scales} with complementary spatial mesh at most
$2^{-k}$.  Its $\delta_n$-diameter tends to zero.  Diagonally
extract all these countably many skeletons, so that their
$j_n$-images converge respectively to points $y_k\in Y$.

Fix $k$.  For any $(y,z),(y',z')\in R$ with $y,y'\ne y_k$,
Lemma~\ref{comp:avoid-separator} joins representing vertices in
the complement of the skeleton.  They therefore lie in the same
spatial complementary component.  Passing to the limit gives
\begin{equation}
 \label{comp:one-cell}
 \rho(z,z')\le2^{-k}.
\end{equation}
This reasoning also applies to two representations of the same
point $y\ne y_k$.

Pass to a subsequence of $k$ with $y_k\to y_0$.  Any two fixed
points of $Y\setminus\{y_0\}$ differ from $y_k$ for all sufficiently
large $k$.  Equation~\eqref{comp:one-cell} shows that all their
spatial representations agree with a single point $z_0$.  There
is at least one such point of $Y$ because $Y$ is nontrivial.  The
surjectivity of $R$ onto $\Chat$ implies
$(y_0,z)\in R$ for every $z\ne z_0$; closure gives it also for
$z=z_0$.  Surjectivity onto $Y$ and closure likewise give
$(y,z_0)\in R$ for every $y\in Y$.  This proves
\eqref{comp:cross-relation}.  Its exceptional row and column are
unique because both factors contain more than one point.

For $z\ne z_0$, the only $Y$-point related to $z$ is $y_0$, so
the only third coordinate above $(z,z)$ in
\eqref{comp:relation-graph} is zero.  At $z_0$ the entire space $Y$
is related, and a pair realizing its positive diameter supplies
\eqref{comp:positive-diagonal}.  The criterion is measurable:
for a closed spatial set $C$, its exceptional point lies in $C$
if and only if, for some integer $m\ge1$, the compact graph hits
the closed set $\{(z,z,r):z\in C,\ r\ge1/m\}$.  These are Borel
conditions in the compact-set topology.  Thus $z_0$ can be retained
without reference to an exploration root or a choice of GH
realization.
\end{proof}

\begin{lemma}[The concentration point avoids the root]
\label{comp:root-avoidance}
In the setting of Proposition~\ref{comp:concentration}, almost
surely $z_0\ne p_*$.
\end{lemma}

\begin{proof}
Keep the original canonical coordinate for this argument.  The
finite pair $(\mu_n,K_n(a_n))$ is determined by the unrooted
decorated map and the three normalization samples.  Conditional
on those data before forgetting incidences, the initial vertex
of a uniform oriented root edge has law
\[
 \nu_n^{\mathrm{root}}
   =\sum_{v\in V_n}\frac{\deg(v)}{2n}\,\delta_{z_n(v)}.
\]
Lemma~\ref{intro:root-flag} couples this distribution to $\mu_n$
at spherical error at most the maximum flag diameter.  Consequently,
for bounded continuous $g$ on $\Chat$ and bounded continuous $F$
of the root-free pair,
\begin{align*}
 &\EE\bigl[F(\mu_n,K_n(a_n))g(z_n(v_*))\bigr]\
 &\hspace{12mm}-
 \EE\left[F(\mu_n,K_n(a_n))\int g\,\dd\mu_n\right]
 \longrightarrow0.
\end{align*}
The difference is bounded by $\|F\|_\infty$ times the modulus of
continuity of $g$ at the maximum flag diameter, which tends to
zero in probability and is bounded.

Pass to the joint limit and use bounded continuous tests followed
by the monotone-class theorem.  Conditional on
$(\mu,\mathcal K(R))$, the point $p_*$ has distribution $\mu$.
The point $z_0$ is measurable from $\mathcal K(R)$ and $\mu$ has
no atoms by Corollary~\ref{conf:mesh-area}.  Hence
\[
 \PP[p_*=z_0]
   =\EE\bigl[\mu(\{z_0\})\bigr]=0.
\]
\end{proof}

We now return to the original oriented peanosphere coordinate.
Uniform convergence of $f_n^{-1}$ gives the geometric relation
\begin{equation}
 \label{comp:oriented-relation}
 R^{\mathrm o}=(\id_Y\times f^{-1})(R)
  =\lim_n\{(j_n(v),x_n(v)):v\in V_n\}.
\end{equation}
In the zero-scale case it is the cross relation with spatial center
$\zeta_0=f^{-1}(z_0)$ and metric center $y_0$.
Lemma~\ref{comp:root-avoidance} implies
$\zeta_0\ne p_*^{\mathrm o}=\infty$.  This pullback is an identity
of locations on each sample; the field remains $h^{\mathrm o}$ with
its original oriented law.

The local-kernel construction retains first-exit observations in
this original coordinate.  For nested spatial disks
$U\Subset V\Subset W$, define
\begin{equation}
 \label{comp:escape-readout}
 E_n^{\mathrm o}(U,V)=
 \max_{\substack{u\in V_n\\x_n(u)\in U}}
 \delta_n\bigl(u,\{v\in V_n:x_n(v)\notin V\}\bigr).
\end{equation}
We use a maximum of zero when the inner set is empty.  The
outer set is nonempty for large $n$ in all tests below.  Distances
in \eqref{comp:escape-readout} equal the minimum cost up to the
first exit from $V$.  With full stars protected in $W$, the
readout uses only that protected graph: the first exiting edge
and its endpoint remain in $W$ once the flag mesh is small.
Lemma~\ref{ker:readouts} retains these quantities, with strict
inner/outer domain sandwiches, in $[0,\infty]$ for the present
deterministic units $a_n$.  We write $E^{\mathrm o}(U,V)$ for such a
retained limit.  All applications use strict containment and strict cost
inequalities, so the sandwiches may be fixed before extraction.

\begin{proposition}[Excluding concentration by two escape tests]
\label{comp:no-concentration}
The alternative $\alpha=0$ is impossible.
\end{proposition}

\begin{proof}
Suppose that the deterministic subsequence in
\eqref{comp:alpha} has $\alpha=0$.  The cross relation
\eqref{comp:cross-relation} then holds almost surely for the
finite-sphere extraction.  We will use an orientation-preserving
reembedding by fresh area marks and produce two separated positive
escape readouts in that reembedded sphere law.  The cross relation
is preserved by the reembedding, so it must hold in the same
experiment.

\paragraph{Finding one local witness.}
The original oriented root is already at infinity, and
$\zeta_0\ne\infty$ by \eqref{comp:oriented-relation}.  Sample two
fresh points $s^1,s^2$ independently from $\mu^{\mathrm o}$,
conditionally on the existing continuum input and output arrays.
Use fresh sampling randomness: their conditional distribution depends
only on the real surface.  They are finite, distinct and different
from $\zeta_0$ almost surely.  Form the orientation-preserving affine
map
\[
 g(z)=\frac{z-s^1}{s^2-s^1}.
\]
This map is defined in the original oriented chart, without using $f$.
Put $\widehat x_n=g\circ x_n$, $\widehat\zeta_0=g(\zeta_0)$,
$\widehat\mu=g_*\mu^{\mathrm o}$, and let $\widehat h$ be the usual
orientation-preserving coordinate transform of $h^{\mathrm o}$ by
$g$.  The geometric relation
$\widehat R=(\id_Y\times g)(R^{\mathrm o})$ is still a cross, now
at $\widehat\zeta_0$ and $y_0$.

We check the law used for this reembedding.  We call it the
\emph{oriented sphere law with fresh area marks}, distinguishing it
from the original coordinate $x_n$ and from the possibly reflected
canonical coordinate.  The real field $\widehat h$ has the ordinary
root-and-two-area-mark sphere law of
Proposition~\ref{enc:sphere-representation}.  The map $g$ depends
only on the real surface and the two fresh sampling variables.
Conditional on those variables, the original imaginary field has
its unchanged independent whole-plane law.  An affine coordinate
change preserves that law with its angular convention: translation
and dilation use the stated symmetries, and rotation adds a constant
to the uniform phase modulo its angular period.  Thus the transformed
imaginary input retains that same law independently of the transformed
real surface.  The new marks are auxiliary continuum variables
sampled independently of the output arrays given the original
continuum input.  Theorem~\ref{ker:local-kernels} therefore retains
the same intrinsic observation kernels in this reembedding.  All
field-density applications below are to $\widehat h$ with this
oriented law.  No reflected field is used.

Let $\widehat E_n$ and $\widehat E$ denote
\eqref{comp:escape-readout} and its retained limit with
$\widehat x_n$ in place of $x_n$.  These still use the same graph
and the same deterministic distance units $a_n$.

Choose disks from a countable ordinary atlas such that
\begin{equation}
 \label{comp:escape-geometries}
 \widehat\zeta_0\in U\Subset V\Subset W,\qquad
 \widehat\mu(\overline W)<\tfrac13,
\end{equation}
and the closed buffer $\overline W$ avoids the marks.  Such choices
exist by diffuseness at $\widehat\zeta_0$ and the preceding mark avoidance.
Since $Y$ is not a point, its maximal distance from $y_0$ is
positive.  Countable selection therefore gives fixed disks and
a fixed $e>0$ for which \eqref{comp:escape-geometries} holds and
some $y\in Y$ has $d_Y(y,y_0)>4e$, on an event of positive
probability.

Represent this point $y$ by vertices $u_n$.  Their spatial
locations in $\widehat x_n$ tend to $\widehat\zeta_0$, so they are
eventually in $U$.  Uniformly
over vertices whose spatial locations are outside $V$, their
$Y$-images tend to $y_0$.  Otherwise a subsequence would produce
a pair of $\widehat R$ with spatial entry outside $V$ and metric
entry different from $y_0$, contradicting its cross form.
Equation~\eqref{comp:approximation} now shows that
$\widehat E_n(U,V)>3e$ for large $n$ on this event.
Using slightly nested disks and the readout sandwiches gives
\begin{equation}
 \label{comp:positive-escape}
 \PP\left[\widehat E(U,V)>2e,\
          \widehat\mu(\overline W)<\tfrac13\right]>0.
\end{equation}

\paragraph{Placing two witnesses in the same oriented sphere law.}
The fixed disks in this paragraph are in the oriented
$\widehat x_n$-coordinate.  The readout in
\eqref{comp:positive-escape} is a local observation
in the sense of Lemma~\ref{ker:readouts}.  First restrict to a
positive-probability subevent on which all imaginary-order data
needed by the test are determined in one finite enlarged
imaginary buffer; Lemma~\ref{ker:finite-order} allows this
restriction.  The domination part of
Lemma~\ref{ker:sphere-density} transfers the positive event to an
ordinary-field experiment with the same local observation
kernel.  No distance or clock normalization is changed.

Translate two copies of this experiment into disks
$U_i\Subset V_i\Subset W_i$, $i=1,2$, whose closed real buffers
and enlarged imaginary buffers are mutually disjoint and avoid
the marks.  Use independent restrictions and their observation
kernels.  Each copy has a positive probability of satisfying the
translated version of \eqref{comp:positive-escape}; hence the
intersection has positive probability.  Its two real-buffer
masses sum to less than $2/3$.

The positivity part of Lemma~\ref{ker:sphere-density}, together
with the product assertion of Theorem~\ref{ker:local-kernels},
puts this event in the actual joint limiting law of the two
finite-sphere readouts.  Proposition~\ref{ker:field-transfer} applies
with translations, whose derivatives are one, so both readouts
still use the original deterministic units $a_n$.  We conclude
that, in the oriented sphere law with fresh area marks,
\begin{equation}
 \label{comp:two-escapes}
 \PP\bigl[\widehat E(U_1,V_1)>2e,\
          \widehat E(U_2,V_2)>2e\bigr]>0.
\end{equation}
This conclusion uses the joint kernel theorem for the two
readouts.  It does not follow from their individual marginals.

\paragraph{The cheap exit contradiction.}
On every sample satisfying the cross relation for $\widehat R$,
at least one of the disjoint closed buffers $\overline W_i$
misses $\widehat\zeta_0$.  Fix
such an index $i$.  Uniformly over vertices in $U_i$, their
$Y$-images tend to $y_0$.  Choose any fixed spatial point of
$W_i\setminus\overline V_i$ and vertices tending to it; their
$Y$-images also tend to $y_0$.  Therefore the maximum distance
from any inner vertex to this exterior target tends to zero,
by \eqref{comp:approximation}.  A shortest graph path to the
target has a first exit from $V_i$ and its initial segment costs
no more than the whole path.  Thus
$\widehat E_n(U_i,V_i)\to0$, including all fixed readout sandwiches,
and $\widehat E(U_i,V_i)=0$.  This contradicts
\eqref{comp:two-escapes}.  Both the two positive readouts and the
cross relation were assertions in the oriented sphere law with
fresh area marks.  Since the affine reembedding preserves the
cross relation sample by sample, the contradiction also excludes
$\alpha=0$ for the canonical relation $R$.
\end{proof}

\subsection{Deterministic normalization and the full sequence}

The final identification first fixes the numerical scale and then
removes the common orientation choice.  The latter uses the following
elementary observation, which also applies to the complete tuple in
Section~\ref{joint:section}.

\begin{lemma}[One common orientation ambiguity]\label{joint:reflection}
Let $E$ be a measurable space with a measurable involution
$\mathfrak r$.  Suppose random elements $X,Y$ satisfy
$X\in\{Y,\mathfrak rY\}$ almost surely, and both $\Law(X)$ and
$\Law(Y)$ are invariant under $\mathfrak r$.  Then
$\Law(X)=\Law(Y)$.  The choice between $Y$ and $\mathfrak rY$ may
depend on the entire coupled random object.
\end{lemma}

\begin{proof}
For any bounded measurable $F:E\to\R$, the unordered pairs
$\{X,\mathfrak rX\}$ and $\{Y,\mathfrak rY\}$ coincide.  Hence
\[
 F(X)+F(\mathfrak rX)=F(Y)+F(\mathfrak rY)
\]
almost surely.  Taking expectations and using both invariances gives
$2\EE F(X)=2\EE F(Y)$.
\end{proof}

\begin{lemma}[Identifying the scalar]
\label{comp:scalar}
If $X$ is a finite, strictly positive random variable and $c>0$
is deterministic, then $cX$ and $X$ have the same law only if
$c=1$.
\end{lemma}

\begin{proof}
If $c>1$, equality of laws implies, for every finite $t>0$,
\[
 \PP[X\le t]=\PP[X\le t/c^k]\longrightarrow\PP[X=0]=0.
\]
This contradicts the almost sure finiteness of $X$.  If $c<1$,
apply the same argument to $1/c$.
\end{proof}

\begin{proof}[Proof of Theorem~\ref{comp:metric-graph}]
Start with an arbitrary deterministic subsequence of positive
integers.  The preceding sections and
Theorem~\ref{intro:intrinsic-input} give a joint further extraction
of the reference data, local readouts and intrinsic spaces.
Their marginals are tight; adjoining the compactified local
arrays preserves tightness.  The full distance graphs are tight
as well: their first two coordinates lie in a fixed compact
sphere, and their third coordinates are bounded by the tight
intrinsic diameters.  Take a Skorohod coupling for these retained
data, and then apply Lemma~\ref{comp:relation} pathwise.

Lemma~\ref{comp:scales} allows the deterministic subsequence to
be chosen so that $a_nB(n)\to\alpha$.  Proposition
\ref{comp:no-concentration} gives $\alpha>0$, and Proposition
\ref{comp:positive-fibers} identifies the location relation as a
homeomorphism satisfying
\[
 (Y,d_Y)\cong(\Chat,\alpha\lambda D).
\]
Both $\alpha$ and $\lambda$ are deterministic: the first is a
limit of numerical normalizations, and the second is the
law-wide rigid comparison constant in
Theorem~\ref{met:normalizers}.  By
\eqref{comp:two-coordinates}, $f$ is an isometry from
$(\Chat,D^{\mathrm o})$ to $(\Chat,D)$ on each sample, regardless
of its orientation or dependence on the data.  Therefore the
unmarked space $(\Chat,D)$ has the ordinary unit-area marginal
of $(\Chat,D^{\mathrm o})$.  This is the same marginal as
$(Y,d_Y)$ by Theorem~\ref{intro:intrinsic-input}.  Their diameters
are finite and strictly positive.  Lemma~\ref{comp:scalar}
gives $\alpha\lambda=1$; no assertion about the law of an
adaptively reembedded field is used here.

Equation~\eqref{comp:relation-graph} is now exactly
\eqref{comp:target-graph}.  In particular, it identifies every
vertex pair in the Hausdorff limit, rather than only a countable
or area-typical set of pairs.  The area convergence is already
joint in Corollary~\ref{conf:mesh-area}.  The resulting canonical
pair is
\[
 X=(\mu,\mathcal G_D)
   \in\{X^+,\mathfrak r X^+\},\qquad
 X^+=(\mu^+,\mathcal G_{D^+}),
\]
where reflection acts simultaneously on the measure and both
spatial coordinates of the distance graph.  We now remove this
possible adaptive sign at the level of the pair's law.

The law of each discrete pair $(\mu_n,K_n(a_n))$ is reflection
invariant: reflection is a weight-preserving bijection of the
oriented finite map ensemble, preserves graph distances and
flag areas, and conjugates the normalized coordinate because
$0,1,\infty$ are fixed by conjugation.  Reflection is continuous
in the weak-measure and compact-graph topologies, so $\Law(X)$
is invariant as well.  The law of $X^+$ is reflection invariant
by the reflection symmetry of the ordinary quantum sphere,
its area and metric, and its three independently sampled marks.
Lemma~\ref{joint:reflection}, applied to this pair, gives
$\Law(X)=\Law(X^+)$ even though the choice between $f=T$ and
$f=\mathfrak r\circ T$ may depend on every retained observation.

Every deterministic subsequence therefore has a further
subsequence with the prescribed pair law.  The subsequence
criterion proves the asserted full-sequence convergence.
The metric identification itself was made in the joint
extraction, so it remains \eqref{comp:target-graph} with any
additional retained observations; the common orientation sign
is kept with those observations.  The normalization $a_n\to0$
is the one from the companion intrinsic theorem.
\end{proof}
